% Options for packages loaded elsewhere
\PassOptionsToPackage{unicode}{hyperref}
\PassOptionsToPackage{hyphens}{url}
\PassOptionsToPackage{dvipsnames,svgnames,x11names}{xcolor}
\documentclass[
  10pt,
]{article}
\usepackage{xcolor}
\usepackage[margin=1in]{geometry}
\usepackage{amsmath,amssymb}
\setcounter{secnumdepth}{-\maxdimen} % remove section numbering
\usepackage{iftex}
\ifPDFTeX
  \usepackage[T1]{fontenc}
  \usepackage[utf8]{inputenc}
  \usepackage{textcomp} % provide euro and other symbols
\else % if luatex or xetex
  \usepackage{unicode-math} % this also loads fontspec
  \defaultfontfeatures{Scale=MatchLowercase}
  \defaultfontfeatures[\rmfamily]{Ligatures=TeX,Scale=1}
\fi
\usepackage{lmodern}
\ifPDFTeX\else
  % xetex/luatex font selection
\fi
% Use upquote if available, for straight quotes in verbatim environments
\IfFileExists{upquote.sty}{\usepackage{upquote}}{}
\IfFileExists{microtype.sty}{% use microtype if available
  \usepackage[]{microtype}
  \UseMicrotypeSet[protrusion]{basicmath} % disable protrusion for tt fonts
}{}
\makeatletter
\@ifundefined{KOMAClassName}{% if non-KOMA class
  \IfFileExists{parskip.sty}{%
    \usepackage{parskip}
  }{% else
    \setlength{\parindent}{0pt}
    \setlength{\parskip}{6pt plus 2pt minus 1pt}}
}{% if KOMA class
  \KOMAoptions{parskip=half}}
\makeatother
\usepackage{longtable,booktabs,array}
\usepackage{needspace,placeins}
\usepackage{caption}
\captionsetup[table]{skip=6pt}
\newcounter{none} % for unnumbered tables
\usepackage{calc} % for calculating minipage widths
% Correct order of tables after \paragraph or \subparagraph
\usepackage{etoolbox}
\makeatletter
\patchcmd\longtable{\par}{\if@noskipsec\mbox{}\fi\par}{}{}
\makeatother
% Allow footnotes in longtable head/foot
\IfFileExists{footnotehyper.sty}{\usepackage{footnotehyper}}{\usepackage{footnote}}
\makesavenoteenv{longtable}
\usepackage{graphicx}
\makeatletter
\newsavebox\pandoc@box
\newcommand*\pandocbounded[1]{% scales image to fit in text height/width
  \sbox\pandoc@box{#1}%
  \Gscale@div\@tempa{\textheight}{\dimexpr\ht\pandoc@box+\dp\pandoc@box\relax}%
  \Gscale@div\@tempb{\linewidth}{\wd\pandoc@box}%
  \ifdim\@tempb\p@<\@tempa\p@\let\@tempa\@tempb\fi% select the smaller of both
  \ifdim\@tempa\p@<\p@\scalebox{\@tempa}{\usebox\pandoc@box}%
  \else\usebox{\pandoc@box}%
  \fi%
}
% Set default figure placement to htbp
\def\fps@figure{htbp}
\makeatother
\setlength{\emergencystretch}{3em} % prevent overfull lines
\providecommand{\tightlist}{%
  \setlength{\itemsep}{0pt}\setlength{\parskip}{0pt}}
\usepackage{bookmark}
\IfFileExists{xurl.sty}{\usepackage{xurl}}{} % add URL line breaks if available
\urlstyle{same}
% fallback for those not using the hyperref driver hyperxmp:
\makeatletter
\@ifundefined{xmpquote}{\newcommand{\xmpquote}[1]{#1}}{}
\makeatother
\hypersetup{
  pdfauthor={Dylan Couzon, Evgeniya Sukhodolskaya, Kumar Shivendu, Andrey Vasnetsov},
  pdftitle={Constella: Lower-Cost Queries over a Fixed Document Index},
  colorlinks=true,
  linkcolor={Maroon},
  filecolor={Maroon},
  citecolor={Blue},
  urlcolor={Blue},
  pdfcreator={LaTeX via pandoc}}

\title{Constella: Lower-Cost Queries over a Fixed Document Index}
\author{\begin{tabular}{cc}
\begin{tabular}{c}Dylan Couzon\\
{\small\href{mailto:dylan.couzon@qdrant.com}{\nolinkurl{dylan.couzon@qdrant.com}}}\end{tabular} & \begin{tabular}{c}Evgeniya Sukhodolskaya\\
{\small\href{mailto:jenny@qdrant.com}{\nolinkurl{jenny@qdrant.com}}}\end{tabular}\\[0.8em]
\begin{tabular}{c}Kumar Shivendu\\
{\small\href{mailto:kumar.shivendu@qdrant.com}{\nolinkurl{kumar.shivendu@qdrant.com}}}\end{tabular} & \begin{tabular}{c}Andrey Vasnetsov\\
{\small\href{mailto:andrey@qdrant.com}{\nolinkurl{andrey@qdrant.com}}}\end{tabular}
\end{tabular}\\[0.6em]Qdrant}
\date{October 2026}

\begin{document}
\maketitle

\subsection{Abstract}\label{abstract}

Can we reduce query computation while keeping an existing dense document
index? Constella provides three compatible query encoders for Stella, a
pretrained 400-million-parameter embedding model: the original encoder,
Nano, a 34.5-million-parameter transformer, and Zero, a token lookup
table. Across 15 BEIR datasets, Nano and Zero retain 90.5\% and 81.4\%
of Stella's exact-search score. On a laptop CPU, their median warmed
encoding times are 2.25 ms and 0.044 ms, versus Stella's 31.6 ms.

We study the quality and systems tradeoffs behind these encoding
savings. Zero requires more approximate-search effort on the two
measured collections; on a million-passage diagnostic, a roughly 62-fold
encoding advantage over Nano becomes 2.2-fold after retrieval, at
different absolute quality levels. A 26-checkpoint table-fitting study
shows why choosing a teacher by its own retrieval score can select a
poor student. Earlier experiments show that fusion-operator ordering
changes with candidate depth and that checkpoint alignment does not
alone certify serving compatibility. Simple routing and vector blending
do not deliver the gains their apparent headroom suggests. The
contribution is an empirical study of cheaper query paths over fixed
document vectors, with findings spanning construction, retrieval, and
deployment, and their limits kept visible.

\subsection{1. Query computation over fixed document
vectors}\label{query-computation-over-fixed-document-vectors}

A large embedding model may be affordable when documents are indexed and
costly on every query. Switching to an unrelated small model usually
means encoding the documents again: the same vector width does not make
two models' coordinates compatible. \textbf{Query-side distillation}
offers another option. A smaller encoder learns to produce queries in
the original model's document space, allowing the stored document
vectors to stay in place.

This paper investigates the quality and retrieval cost of that option.
We built \textbf{Constella} around \textbf{Stella}, the pretrained
\texttt{stella\_en\_400M\_v5} embedding model {[}Stella{]}. Its document
vectors are fixed. \textbf{Nano} is a small transformer aligned to those
vectors. \textbf{Zero} looks up and pools learned token vectors, without
running a transformer. The original Stella query encoder remains
available. An application can choose any of these aligned encoders for a
request and search the same collection.

Index reuse and lookup-table queries have precedents. QED, EmbedDistill,
LEAF, and NanoVDR study compact query encoders; pyNIFE aligns a static
table to a frozen teacher and proposes static/contextual switching
{[}QED, EmbedDistill, LEAF, NanoVDR, pyNIFE{]}. Model2Vec and static
sentence embeddings supply other table-based representations;
LightRetriever learns a lookup query side together with its document
side {[}Model2Vec, StaticEmb, LightRetriever{]}. We claim empirical
findings and usable artifacts, not a new architecture or priority for
switching.

The useful choices extend beyond encoder size. Which teacher should a
cheap encoder imitate? Does the existing approximate-search
configuration still work? Would a lexical branch buy more quality than
additional query computation? Does the exported model preserve the
vectors that were evaluated?

\Needspace{8\baselineskip}
Our results establish four findings:

\begin{itemize}
\tightlist
\item
  \textbf{Several query budgets can share the same index.} Nano and Zero
  offer large encoding savings with measurable relevance loss; keeping
  Stella's index is their common constraint.
\item
  \textbf{Teacher quality does not rank table quality reliably under the
  tested recipe.} A direct table screen is more informative than teacher
  scores or vector agreement.
\item
  \textbf{Search work can consume much of the encoding saving.}
  Additional candidate exploration reduces Zero's advantage; a
  controlled query-precision test identifies recoverable candidate
  information.
\item
  \textbf{Encoder quality alone does not determine the deployed result.}
  Fusion depth changes operator ordering, simple adaptive policies
  underperform their oracle, and serving wrappers can violate
  compatibility.
\end{itemize}

The following sections present the experiments and their implications.
``Registered'' means the method was fixed before observing the result;
``exploratory'' identifies later analyses. Detailed methods and the
original registered evaluations are in the appendices; the companion
evidence package preserves the broader research history.

\subsection{2. What the three query paths
offer}\label{what-the-three-query-paths-offer}

\subsubsection{2.1 One compatibility
contract}\label{one-compatibility-contract}

Stella produces normalized 1024-dimensional document vectors. Zero and
Nano are trained against that particular space. Zero has 30,522 learned
token rows, exported as int8. It pools token vectors with square-root
count weights, so repeated tokens gain influence more slowly than simple
counting, then normalizes. Nano has 34,540,672 parameters: it starts
from bge-small, a pretrained small English embedding model, combines
three of its transformer layers, projects their outputs to 1024
dimensions, averages token outputs, and learns to match Stella query
vectors using squared L2 loss.

Here, \textbf{hot swapping} means selecting an available, aligned
encoder without re-encoding documents or building another dense index.
We verified all three query paths against the same populated Qdrant
collections, with no updates or rebuild between encoder batches. This
demonstrates compatibility; model-loading transitions and concurrent
failover were not timed. The prompt, tokenizer, pooling, normalization,
and model revision remain part of the contract.

\subsubsection{2.2 Quality and encoding
cost}\label{quality-and-encoding-cost}

We use \textbf{nDCG@10}, which rewards relevant documents near the top
of the first ten results. The BEIR-15 macro average gives equal weight
to each of 15 retrieval datasets {[}BEIR{]}. Exact search scores every
document vector, separating encoder quality from approximate-search
errors. \textbf{Retention} is the ratio of a student's macro score to
Stella's, not a percentage of queries answered correctly.

\Needspace{16\baselineskip}
\textbf{Table 1.} Three query paths over identical Stella document
vectors. Quality is a descriptive BEIR-15 aggregate. Encoding medians
follow a registered protocol: 100 real 5--12-word queries, batch one,
four CPU threads, Apple M5 Pro, three fresh processes. Transformers use
ONNX Runtime; Zero uses NumPy. Asset sizes include model and tokenizer
files.

{\def\LTcaptype{none} % do not increment counter
\begin{longtable}[]{@{}lrrrr@{}}
\toprule\noalign{}
Query path & Exact nDCG@10 & Retention & Encode p50 & Assets \\
\midrule\noalign{}
\endhead
\bottomrule\noalign{}
\endlastfoot
Stella & 0.5614 & 100\% & 31.6 ms & 1669.6 MiB \\
Nano & 0.5081 & 90.5\% & 2.25 ms & 132.3 MiB \\
Zero & 0.4572 & 81.4\% & 0.044 ms & 90.1 MiB \\
\end{longtable}
}

Nano keeps about nine tenths of Stella's score at one fourteenth of its
encoding time. Zero keeps about four fifths at one seven-hundredth.
These are different relevance budgets, not equal-quality speedups.
Zero's warmed encoding is small, but tokenization, asset loading,
memory, and retrieval remain. Its measured load time is 215 ms; the
first query takes 0.208 ms.

Keeping the index is the reason to consider this family. In a new-index
decision, alternatives deserve their own evaluation. Nano's BEIR-15
score is below the selected bge-small and LEAF references, which use
their own document representations (Appendix A). BM25 and inference-free
sparse retrieval also offer cheap queries. Constella is not a
comprehensive small-model ranking.

\subsection{3. Encoding savings and approximate-search
cost}\label{encoding-savings-and-approximate-search-cost}

\subsubsection{3.1 Where the encoding saving
goes}\label{where-the-encoding-saving-goes}

Approximate nearest-neighbor search (\textbf{ANN}) explores a limited
set of candidates rather than scoring every document. We measured all
three encoders with Qdrant 1.19.1 on FiQA (57,638 documents) and a
one-million-passage MS MARCO diagnostic. The latter retains every judged
positive and samples the remaining passages, removing most full-corpus
competitors. Its absolute scores and search penalties are not
full-corpus MS MARCO results; it remains validation-only.

The registered sweep varies graph-search effort and compression. HNSW's
\texttt{ef} controls the search exploration budget. With compression,
\textbf{oversampling} widens the candidates reconsidered using original
vectors. Each compression configuration has its own collection; within a
collection, documents and graph remain fixed for all query encoders.

At low search effort (\texttt{ef=16}) on the uncompressed
million-passage collection, ANN loses 16.3\% of Zero's exact-search
nDCG@10, compared with 7.0\% for Nano and 4.4\% for Stella. Raising
search effort repairs much of this loss, but spends time. Earlier
LightRetriever experiments showed a similar lookup-query penalty on
FiQA; the size of the effect varied by dataset. Query-distribution
effects on graph search are established in OOD-DiskANN and RoarGraph
{[}OOD-DiskANN, RoarGraph{]}. These measurements test their practical
consequence for this compatible family.

\Needspace{14\baselineskip}
\textbf{Table 2.} On one shared binary collection, the fastest recorded
settings within 1\% of each encoder's own exact score give:

{\def\LTcaptype{none} % do not increment counter
\begin{longtable}[]{@{}llrr@{}}
\toprule\noalign{}
Query path & \texttt{ef} / oversampling & ANN nDCG@10 & Encode + search
p50 \\
\midrule\noalign{}
\endhead
\bottomrule\noalign{}
\endlastfoot
Zero & 512 / 2× & 0.6116 & 1.11 ms \\
Nano & 256 / 4× & 0.6893 & 2.48 ms \\
Stella & 128 / 1× & 0.7226 & 21.3 ms \\
\end{longtable}
}

For example, suppose an application required a score of at least 0.68 on
this diagnostic and a median encoding-plus-search budget of 3 ms. Among
these three measured settings, Nano meets both. Zero falls below the
relevance floor; Stella exceeds the time budget. If the budget were 2
ms, none of these rows would meet both requirements. These hypothetical
thresholds illustrate the choice, not recommended production targets.

This exploratory restriction of the registered sweep keeps the physical
index fixed. On these queries, Nano encoding is roughly 62 times slower
than Zero; after search, the gap is 2.2-fold. Table 1 uses another query
sample and implies a 51-fold encoding gap. On the uncompressed
collection, Zero and Nano require 3.61 and 3.85 ms at the same
tier-relative target, nearly consuming the advantage. FiQA's best
measured configurations leave a larger 3.8--4.3-fold gap.

\begin{figure}
\centering
\pandocbounded{\includegraphics[keepaspectratio,alt={Measured quality/cost choices from the registered sweep. Curves can use different compression configurations; dotted lines show each encoder's exact score. The shared-binary table above restricts all tiers to one collection. Time sums separately measured per-query encoding and search phases, excluding model loading and application dispatch. The quality levels are different.}]{figures/f4_system.pdf}}
\caption{Measured quality/cost choices from the registered sweep. Curves
can use different compression configurations; dotted lines show each
encoder's exact score. The shared-binary table above restricts all tiers
to one collection. Time sums separately measured per-query encoding and
search phases, excluding model loading and application dispatch. The
quality levels are different.}
\end{figure}

The result separates compatibility from search efficiency. Shared
document vectors do not establish that the original ANN settings remain
sufficient. A target relative to each encoder's own score isolates
additional approximation loss; the absolute relevance requirement still
determines which encoder is usable.

\FloatBarrier
\subsubsection{3.2 Query precision deserves its own
test}\label{query-precision-deserves-its-own-test}

A binary document code stores coordinate signs. The tested engine's
default binary query encoding also keeps only signs, discarding query
magnitudes {[}Qdrant 1.19.1{]}. We held document codes and query vectors
fixed, removed the graph, and exhaustively compared \texttt{sign(q)} dot
\texttt{sign(d)} with normalized \texttt{q} dot the same
\texttt{sign(d)}. The target remained each encoder's original-vector
exact top ten.

At a 40-candidate budget, preserving query magnitudes raises Zero's
original-neighbor coverage from 82.2\% to 93.0\% on FiQA and from 90.1\%
to 96.8\% on the million-passage diagnostic. All three encoders improve.
Zero starts with more misses and does not repair a larger fraction of
them than the other tiers. The control shows that its cheap vectors
still contain useful magnitude information; it does not measure a native
latency or relevance gain. Query-aware binary scoring already has
precedents {[}QA-Cos{]}. Appendix B preserves the controls and the
weaker native replication.

Query precision is therefore a separate budget from encoder compute and
document compression. A native comparison at matched quality and
measured cost remains necessary to establish a faster production
setting.

\FloatBarrier
\subsection{4. Lexical fusion and adaptive query
computation}\label{lexical-fusion-and-adaptive-query-computation}

\subsubsection{4.1 Lexical fusion is a practical
alternative}\label{lexical-fusion-is-a-practical-alternative}

Combining Zero with BM25 raises its BEIR-15 macro score from 0.4572 to
0.4933, or 87.9\% of Stella's score. Nano plus BM25 reaches 0.5110.
These use distribution-based score fusion (\textbf{DBSF}), which
normalizes and merges dense and lexical candidate scores, with 100
candidates per branch. Fusion adds a sparse index and retrieval path. On
the million-passage diagnostic, a separate uncompressed
dense-plus-sparse collection takes 4.71 ms including Zero encoding,
versus 1.11 ms for the dense binary example. These are different
configurations, not a pure fusion overhead measurement.

Fusion's quality also varies by workload: HotpotQA Zero improves from
0.6127 to 0.7015, while MS MARCO Nano falls from 0.4063 to 0.3699. A
later internal specialization attempt improved dense development results
but failed its fused serving gate. Evaluate the route that will actually
serve queries.

Earlier fusion research supplies a less obvious choice:
\textbf{candidate depth changes the observed operator ordering}. On four
development components, DBSF slightly exceeds fitted convex score fusion
at ten candidates but trails it at a thousand. Convex fusion rescales
each branch's scores and mixes them with the fixed development-fitted
weight 0.8; DBSF uses the branch's score distribution without fitting
that weight. The lexical branch uses bm25s with Lucene settings. Other
lexical implementations need their own comparison because their score
distributions enter DBSF normalization.

\Needspace{13\baselineskip}
\textbf{Table 3.} Fusion scores on the same development components, at
different candidate depths.

{\def\LTcaptype{none} % do not increment counter
\begin{longtable}[]{@{}lrrr@{}}
\toprule\noalign{}
Candidates per branch & Convex fusion & DBSF & DBSF − convex \\
\midrule\noalign{}
\endhead
\bottomrule\noalign{}
\endlastfoot
10 & 0.5482 & 0.5517 & +0.0035 \\
50 & 0.5578 & 0.5558 & −0.0020 \\
100 & 0.5637 & 0.5574 & −0.0063 \\
1000 & 0.5727 & 0.5580 & −0.0146 \\
\end{longtable}
}

This development curve is descriptive. It does not establish equivalence
at shallow depth or isolate a causal scoring mechanism. Our eventual
DBSF@100 recommendation was an implementability decision after earlier
evaluation; the original registered convex-fusion result remains in
Appendix A. The lesson is to compare operators at the candidate budget
and lexical implementation being deployed, including any self-document
exclusion.

\subsubsection{4.2 Routing and blending still need
evidence}\label{routing-and-blending-still-need-evidence}

Compatible encoders permit a per-query choice, but an oracle using
relevance labels is not a deployable policy. On 12 nonreserved BEIR
datasets, a Zero/Nano oracle reaches 0.5473 versus always-Nano's 0.5161.
It matches always-Nano by sending the most beneficial 15\% of each
dataset's queries to Nano. The simple policies we tested capture little
of that opportunity.

Subwords per word, query length, and the pooled vector norm are weak
selectors. Zero's post-search score margin recovers about one quarter of
the oracle gain over random routing on six datasets; escalated queries
require another encoding and search. A cheap vector blend also
disappoints: a development-selected 30\% Zero contribution improves Nano
locally but lowers its six-set score by 0.0051 (95\% paired query
interval −0.0092 to −0.0011).

These experiments distinguish available complementarity from realized
benefit. Static encoder choices and lexical fusion have measured
outcomes here; routing adds repeated searches, and the selected blend
loses quality outside development. Better policies remain possible,
consistent with work on retrieval-strategy selection and query
performance prediction {[}Arabzadeh et al., QPP{]}.

\subsection{5. Teacher selection and construction
cost}\label{teacher-selection-and-construction-cost}

\subsubsection{5.1 Teacher quality and table
quality}\label{teacher-quality-and-table-quality}

An existing index fixes the teacher. Before building an index, a team
can also choose which teacher its cheap query encoder will imitate.
Earlier project screens reversed a teacher choice after evaluating its
table. We followed that observation with a larger comparison: one
\textbf{closed-form table recipe} across 26 checkpoints.

The probe fits token vectors using ridge regression, reconstructing
teacher query vectors from token counts with a penalty against large
changes from the starting rows. It uses 337,981 screened fit queries.
The penalty is selected separately for each teacher on two development
forums, then frozen before evaluation on six public BEIR datasets. Ten
checkpoints were registered before observation; 16 were exploratory
additions. One additional checkpoint failed the solver's convergence
requirement and is excluded. All share a BERT WordPiece vocabulary. Each
table searches its own teacher's document vectors: this is a
teacher-plus-index selection study, not a way to swap unrelated teachers
into Stella's index.

\begin{figure}
\centering
\pandocbounded{\includegraphics[keepaspectratio,alt={A direct table test ranks table retrieval much better than the teacher's own score under this recipe. Filled points are the ten registered checkpoints; hollow points are 16 exploratory additions. Spearman correlation measures agreement between rankings. Intervals resample checkpoints, not training runs.}]{figures/f3_towers.pdf}}
\caption{A direct table test ranks table retrieval much better than the
teacher's own score under this recipe. Filled points are the ten
registered checkpoints; hollow points are 16 exploratory additions.
Spearman correlation measures agreement between rankings. Intervals
resample checkpoints, not training runs.}
\end{figure}

\FloatBarrier
The strongest registered teacher, gte-large-en-v1.5, scores 0.5970 on
the six datasets but produces the weakest registered table, at 0.2455.
Stella's teacher scores 0.5745 and its table 0.3974. Selecting the
strongest teacher by that public score would therefore yield a
substantially worse table. This is a hindsight comparison illustrating a
selection failure, not a prospectively tested model-card policy.

Across all 26 checkpoints, teacher/table rank correlation is +0.09 (95\%
interval −0.39 to +0.52). The wide interval does not establish
independence. Development-table/public-table correlation is +0.88 (+0.70
to +0.95), and +0.90 for the registered ten. The development test
selects Stella, the best six-set average table. It is not best on every
dataset; on the four sets without disclosed Stella overlap, ranking
agreement drops to +0.68 and its selected table trails the best
candidate. A representative development workload remains necessary.

Teacher strength is also not the same as ease of imitation. Tables
fitted to two e5 checkpoints match teacher query vectors at mean cosine
0.90 and 0.89, versus 0.78 for Stella's table, yet retrieve worse.
Earlier Zero experiments and a later vocabulary-training screen likewise
reduced training losses without improving ranking. These observations
are consistent with prior work on stronger teachers producing weaker
students {[}Cho and Hariharan 2019, PROD{]}. Our contribution is the
consequential reversal and direct screen under this table recipe, rather
than a general law of distillation.

The direct screen measures the representation whose retrieval is being
selected. Teacher score, imitation fidelity, and optimization loss
measure different outcomes. The probe predicts its own table recipe; it
has not been validated as a selector for the separately trained Zero or
Nano models.

\subsubsection{5.2 Component build costs}\label{component-build-costs}

The published models require no student training to use. Building
another family requires texts, teacher-produced targets, fitting,
export, and retrieval validation. Final Nano optimization took 57.3
hours on one A100 GPU for exactly 199,999,721 example presentations:
about \$95 at the recorded historical rental rate. That is the
optimization component, excluding target preparation, data generation,
recipe search, failed runs, evaluation, and engineering.

Zero's recorded retraining estimate was approximately 20 minutes on an
RTX 3080 with prepared targets; changing teachers was estimated to
require another 8--12 hours of target encoding. Those are historical
estimates rather than measured complete builds. These fitting components
can run on a single GPU; they do not establish a turnkey reproduction
budget or the minimum training needed.

\FloatBarrier
\subsection{6. Serving compatibility}\label{serving-compatibility}

An aligned checkpoint can still produce incompatible queries through the
wrong serving path. In this project, missing Stella's query prompt
changes vector direction; padding to a fixed length corrupts a table
path; a custom Nano bridge passes while native registration initially
uses the wrong pooling. A candidate fp16 document export passes CPU
parity but fails on CUDA, with minimum cosine 0.662 against its
reference. We ship the fp32 path. These are implementation
counterexamples, not claims that fp16 or exports generally fail.

The output contract also includes dtype. A mean-pooling wrapper promotes
float32 embeddings to float64 through mask arithmetic, doubling vector
bytes despite close numerical agreement. Naively narrowing before
normalization can overflow float16 norms. Numerical agreement alone does
not certify the output contract. Our serving qualification includes
values, dtype, boundary lengths, and batch-size changes through the
actual caller and device. Compatibility learning is established {[}BCT,
FCT{]}; these failures concern preservation of the particular alignment
being claimed.

\subsection{7. Limits and conclusion}\label{limits-and-conclusion}

Constella's budgets are measured properties of these trained artifacts.
Zero and Nano use different training pools, including different FEVER
exposure; their comparison does not isolate architecture. Stella also
has disclosed exposure on some benchmark datasets. Appendix A separates
original registered contrasts from later descriptive aggregates and
preserves failed or unresolved tests.

Small internal specialization and patch experiments also encountered
source-object leakage and judgment uncertainty. They did not establish
an improved encoder and remain supporting evidence about evaluation
limits in the companion research history.

The ANN and precision findings repeat across two workloads in one
teacher space, on one laptop and sequential search. The million-passage
subset removes competitors. Query intervals condition on fixed trained
models; checkpoint intervals do not establish an independent-family law.
The teacher screen uses one table recipe and is not validated for
trained transformers. Concurrent switching, all-tier resident memory,
server throughput, and a native precision remedy remain unmeasured.

The practical result is a choice of query budgets without replacing
Stella's document vectors. Making that choice useful requires testing
the student's retrieval, the candidate-search cost, the fused route, and
the actual runtime. The next experiment with the clearest deployment
value is a native query-precision comparison at matched recovery and
relevance, including scoring cost. The current evidence justifies that
experiment; it does not promise its outcome.

\subsection{Appendix A. Evaluation and construction
details}\label{appendix-a.-evaluation-and-construction-details}

\textbf{Serving configuration.} Stella is pinned to revision
\texttt{ffeb2b7e}, with its \texttt{s2p\_query} prompt, fp32 encoding,
normalized output, fp16 stored document vectors, and cosine similarity.
Maximum sequence length is 512. Nano uses bge-small layers 12, 8, and 4,
concatenated to 1152 dimensions before projection; batches mix 75\%
queries and 25\% documents. Zero's training pool includes FEVER-train
and HotpotQA; Nano excludes FEVER. MS MARCO is validation-only in our
distillation, although Nano's pretrained base has prior exposure.
Targets and permitted training sources are recorded in the companion
build records.

\textbf{Partitions and interpretation.} The six-set suite is NFCorpus,
SCIDOCS, SciFact, TREC-COVID, FiQA, and ArguAna. ``Clean-four'' means
the first four have no disclosed Stella overlap; it is not a causal
contamination estimate. Four originally reserved datasets were evaluated
once: FEVER, DBpedia-entity, CQADupStack android, and english. Their
published aggregates enter the broader descriptive evaluation and do not
tune the paper's diagnostics. NDO-3 weights DBpedia 0.5 and the two
forums 0.25 each. Two other forums, physics and programmers, select
table penalties and routing/blend settings. Those development sources
were repeatedly used during earlier model construction.

The external references below were selected project targets, not a
comprehensive or state-of-the-art frontier. They use their own document
spaces. Nano's lower broad score is retained even though it passes some
earlier registered contrasts.

{\def\LTcaptype{none} % do not increment counter
\begin{longtable}[]{@{}
  >{\raggedright\arraybackslash}p{(\linewidth - 6\tabcolsep) * \real{0.2143}}
  >{\raggedright\arraybackslash}p{(\linewidth - 6\tabcolsep) * \real{0.2143}}
  >{\raggedleft\arraybackslash}p{(\linewidth - 6\tabcolsep) * \real{0.2857}}
  >{\raggedleft\arraybackslash}p{(\linewidth - 6\tabcolsep) * \real{0.2857}}@{}}
\toprule\noalign{}
\begin{minipage}[b]{\linewidth}\raggedright
Reference
\end{minipage} & \begin{minipage}[b]{\linewidth}\raggedright
Document representation
\end{minipage} & \begin{minipage}[b]{\linewidth}\raggedleft
BEIR-15 macro nDCG@10
\end{minipage} & \begin{minipage}[b]{\linewidth}\raggedleft
Encode p50
\end{minipage} \\
\midrule\noalign{}
\endhead
\bottomrule\noalign{}
\endlastfoot
bge-small-en-v1.5 & its own dense vectors & 0.5171 & 2.57 ms \\
LEAF & arctic-embed-m dense vectors & 0.5402 & 1.39 ms \\
BM25 & lexical terms & 0.4006 & not reported here \\
\end{longtable}
}

Nano's registered tests use a fixed sequence: each superiority conjunct
is evaluated using a one-sided 2.5\% lower bound at alpha 0.025. The
original partitions and outcomes follow. Held-out intervals and other
reported paired query intervals resample queries with trained models
fixed. They do not include training-seed variation.

{\def\LTcaptype{none} % do not increment counter
\begin{longtable}[]{@{}
  >{\raggedright\arraybackslash}p{(\linewidth - 6\tabcolsep) * \real{0.2308}}
  >{\raggedright\arraybackslash}p{(\linewidth - 6\tabcolsep) * \real{0.2308}}
  >{\raggedleft\arraybackslash}p{(\linewidth - 6\tabcolsep) * \real{0.3077}}
  >{\raggedright\arraybackslash}p{(\linewidth - 6\tabcolsep) * \real{0.2308}}@{}}
\toprule\noalign{}
\begin{minipage}[b]{\linewidth}\raggedright
Original Nano contrast
\end{minipage} & \begin{minipage}[b]{\linewidth}\raggedright
Partition
\end{minipage} & \begin{minipage}[b]{\linewidth}\raggedleft
nDCG@10 difference
\end{minipage} & \begin{minipage}[b]{\linewidth}\raggedright
Outcome
\end{minipage} \\
\midrule\noalign{}
\endhead
\bottomrule\noalign{}
\endlastfoot
Nano − bge-small & clean-four & +0.017648; one-sided lower bound
+0.003674 & established \\
Nano − bge-small & all six & +0.027449 & established \\
Nano − LEAF & all six & +0.016181 & established \\
Nano − LEAF & clean-four & −0.001063 & unresolved \\
Nano − bge-small & held-out NDO-3, dataset weighted & +0.0032
{[}−0.0069, +0.0134{]} & registered descriptive \\
Nano − bge-small & held-out NDO-3, query pooled & −0.0121 {[}−0.0219,
−0.0024{]} & registered descriptive \\
Nano − LEAF & held-out NDO-3 & −0.0389 {[}−0.0488, −0.0290{]} &
registered descriptive \\
\end{longtable}
}

\Needspace{17\baselineskip}
The original Zero tests used a Holm multiplicity rule across the
confirmatory family. The table preserves their raw reporting intervals;
these alone are not the decision rule.

{\def\LTcaptype{none} % do not increment counter
\begin{longtable}[]{@{}
  >{\raggedright\arraybackslash}p{(\linewidth - 6\tabcolsep) * \real{0.2143}}
  >{\raggedleft\arraybackslash}p{(\linewidth - 6\tabcolsep) * \real{0.2857}}
  >{\raggedleft\arraybackslash}p{(\linewidth - 6\tabcolsep) * \real{0.2857}}
  >{\raggedright\arraybackslash}p{(\linewidth - 6\tabcolsep) * \real{0.2143}}@{}}
\toprule\noalign{}
\begin{minipage}[b]{\linewidth}\raggedright
Original Zero contrast, all six
\end{minipage} & \begin{minipage}[b]{\linewidth}\raggedleft
nDCG@10 difference
\end{minipage} & \begin{minipage}[b]{\linewidth}\raggedleft
Raw 95\% interval
\end{minipage} & \begin{minipage}[b]{\linewidth}\raggedright
Registered outcome
\end{minipage} \\
\midrule\noalign{}
\endhead
\bottomrule\noalign{}
\endlastfoot
Zero − LightRetriever dense & −0.0243 & {[}−0.0405, −0.0086{]} &
established below the 0.4583 bar \\
Zero − BM25 & +0.0165 & {[}+0.0017, +0.0311{]} & unresolved; p=.0149
fails Holm .0083 \\
Zero + BM25 convex − OpenSearch & +0.0043 & {[}−0.0063, +0.0151{]} &
unresolved \\
\end{longtable}
}

The registered clean-four sensitivity is descriptive: the corresponding
differences are −0.0443 {[}−0.0675, −0.0212{]}, −0.0311 {[}−0.0517,
−0.0109{]}, and −0.0107 {[}−0.0262, +0.0043{]}. Zero therefore did not
confirmatorily beat BM25. Unresolved superiority does not establish
equivalence. The later DBSF recommendation does not replace the original
confirmatory convex operator (weight 0.8, depth 1000).

\Needspace{20\baselineskip}
The originally reserved four, copied from their published aggregate
receipt:

{\def\LTcaptype{none} % do not increment counter
\begin{longtable}[]{@{}lrrrr@{}}
\toprule\noalign{}
System & FEVER & DBpedia & CQADup android & CQADup english \\
\midrule\noalign{}
\endhead
\bottomrule\noalign{}
\endlastfoot
Nano & 0.6231 & 0.4190 & 0.4774 & 0.4298 \\
Zero & 0.6978 & 0.3900 & 0.4431 & 0.3690 \\
Stella & 0.8207 & 0.4603 & 0.5275 & 0.5183 \\
bge-small & 0.8662 & 0.4003 & 0.4761 & 0.4558 \\
LEAF & 0.8671 & 0.4520 & 0.5260 & 0.4707 \\
BM25 & 0.5036 & 0.3045 & 0.3952 & 0.3481 \\
Zero + BM25 & 0.7559 & 0.4030 & 0.4629 & 0.4016 \\
Nano + BM25 & 0.6879 & 0.4296 & 0.4750 & 0.4323 \\
\end{longtable}
}

\textbf{Teacher-screen limits.} Checkpoint expansion follows the same
selection order and recipe. The fit list was cleaned after an earlier
version had 1.31\% held-out-query overlap. Registered checkpoint
exposure is documented; the exploratory expansion was not individually
audited. A family-resampling sensitivity accompanies checkpoint
intervals. Three vector-error diagnostics do not reliably explain
teacher ranking. Direct table fitting takes roughly 4--7 minutes per
registered configuration on an A100, inferred from output-file intervals
and including query encoding except cached Stella targets; downloads,
data assembly, and index construction are excluded.

\subsection{Appendix B. Search and precision
controls}\label{appendix-b.-search-and-precision-controls}

The ANN sweep indexes every segment before timing, with HNSW
\texttt{m=16}, \texttt{ef\_construct=100}, and an indexing threshold of
1 KB rather than the 10,000 KB default. This avoids plain-scan segments
in the benchmark and is not a deployment recommendation. Exact parity is
checked against the original-vector NumPy reference. The
positive-preserving million-passage sample is deterministic. Timing uses
warmed, sequential queries and separately measured encoding/search
phases; it does not include application dispatch, loading, or concurrent
traffic.

The native binary follow-up alternates options on one fixed collection
per workload, using 192 deterministically selected queries each. At
primary \texttt{ef=64}, binary scoring with rescoring and 1×
oversampling loses 11.7/5.1/2.3 percentage points of original-neighbor
recovery for Zero/Nano/Stella on FiQA. The million-passage replication
loses 3.9/3.1/2.2 points; the extra Zero-versus-Nano interval includes
zero. This weaker replication prevents a general static-model
quantization claim. Native rescoring also changes cross-segment
candidate merging, so it is not a rerank of an identical global pool.

\Needspace{16\baselineskip}
The graph-free precision control uses exactly the same queries, vectors,
and targets. It scores all fixed document sign codes, with global
candidate budgets 10/40/100 and 40 primary. Expected coverage averages
uniform inclusion at tied boundaries. The primary rows are:

{\def\LTcaptype{none} % do not increment counter
\begin{longtable}[]{@{}
  >{\raggedright\arraybackslash}p{(\linewidth - 8\tabcolsep) * \real{0.1667}}
  >{\raggedright\arraybackslash}p{(\linewidth - 8\tabcolsep) * \real{0.1667}}
  >{\raggedleft\arraybackslash}p{(\linewidth - 8\tabcolsep) * \real{0.2222}}
  >{\raggedleft\arraybackslash}p{(\linewidth - 8\tabcolsep) * \real{0.2222}}
  >{\raggedleft\arraybackslash}p{(\linewidth - 8\tabcolsep) * \real{0.2222}}@{}}
\toprule\noalign{}
\begin{minipage}[b]{\linewidth}\raggedright
Workload
\end{minipage} & \begin{minipage}[b]{\linewidth}\raggedright
Query path
\end{minipage} & \begin{minipage}[b]{\linewidth}\raggedleft
Sign query
\end{minipage} & \begin{minipage}[b]{\linewidth}\raggedleft
Full query magnitudes
\end{minipage} & \begin{minipage}[b]{\linewidth}\raggedleft
Gain, percentage points
\end{minipage} \\
\midrule\noalign{}
\endhead
\bottomrule\noalign{}
\endlastfoot
FiQA & Zero & 0.8217 & 0.9297 & +10.80 \\
FiQA & Nano & 0.9137 & 0.9771 & +6.33 \\
FiQA & Stella & 0.9638 & 0.9922 & +2.84 \\
1M diagnostic & Zero & 0.9010 & 0.9677 & +6.67 \\
1M diagnostic & Nano & 0.9484 & 0.9833 & +3.49 \\
1M diagnostic & Stella & 0.9685 & 0.9932 & +2.47 \\
\end{longtable}
}

The paired extra Zero/Nano gains are 4.46 points on FiQA (95\% query
interval 2.53--6.43) and 3.17 on the diagnostic (1.58--4.75). Mean
cosine between full and sign-only queries is about 0.80 for every tier;
the mechanism remains unresolved. Full query values are not native
scalar8 encoding; global candidate counts are not per-segment
oversampling. All precision follow-ups are exploratory. Registered
search and routing methods are distinguished from these controls in the
evidence package.

\subsection{Artifact availability}\label{artifact-availability}

The
\href{https://github.com/Dylancouzon/asymmetric-dual-encoders/tree/m15-whitepaper}{companion
repository} contains inference examples, build configurations, methods,
immutable result receipts, per-dataset tables, claim-to-evidence
mappings, and the full milestone learning catalog. The released
artifacts are
\href{https://huggingface.co/Qdrant/constella-zero}{Constella Zero},
\href{https://huggingface.co/Qdrant/constella-nano}{Constella Nano}, and
the
\href{https://huggingface.co/Qdrant/stella-en-400M-v5-doc-onnx}{Stella
document encoder}. Model hashes and reference-vector checks identify the
evaluated paths. The repository records construction but is not a
standalone trainer for an arbitrary teacher.

\subsection{References}\label{references}

\begin{itemize}
\item
  {[}QED{]} Yuxuan Wang and Hong Lyu. Query Encoder Distillation via
  Embedding Alignment is a Strong Baseline Method to Boost Dense
  Retriever Online Efficiency. SustaiNLP Workshop at ACL, 2023.
  \href{https://arxiv.org/abs/2306.11550}{arXiv:2306.11550}.
\item
  {[}EmbedDistill{]} Seungyeon Kim, Ankit Singh Rawat, Manzil Zaheer, et
  al.~EmbedDistill: A Geometric Knowledge Distillation for Information
  Retrieval. 2023.
  \href{https://arxiv.org/abs/2301.12005}{arXiv:2301.12005}.
\item
  {[}LEAF{]} Robin Vujanic and Thomas Rueckstiess. LEAF: Knowledge
  Distillation of Text Embedding Models with Teacher-Aligned
  Representations. 2025, revised 2026.
  \href{https://arxiv.org/abs/2509.12539}{arXiv:2509.12539}.
\item
  {[}NanoVDR{]} Zhuchenyang Liu, Yao Zhang, and Yu Xiao. NanoVDR:
  Distilling a 2B Vision-Language Retriever into a 70M Text-Only Encoder
  for Visual Document Retrieval. 2026.
  \href{https://arxiv.org/abs/2603.12824}{arXiv:2603.12824}.
\item
  {[}LightRetriever{]} Guangyuan Ma, Yongliang Ma, Xuanrui Gou, Zhenpeng
  Su, Ming Zhou, and Songlin Hu. LightRetriever: A LLM-based Text
  Retrieval Architecture with Extremely Faster Query Inference. ICLR,
  2026. \href{https://arxiv.org/abs/2505.12260}{arXiv:2505.12260}.
\item
  {[}pyNIFE{]} Stephan Tulkens. pyNIFE: Nearly Inference Free Embeddings
  in Python. Software, 2025.
  \href{https://github.com/stephantul/pynife}{Repository}.
\item
  {[}Model2Vec{]} MinishLab. Model2Vec: Fast State-of-the-Art Static
  Embeddings. Software.
  \href{https://github.com/MinishLab/model2vec}{Repository}.
\item
  {[}StaticEmb{]} Tom Aarsen. Train 400x faster Static Embedding Models
  with Sentence Transformers. Hugging Face, 15 January 2025.
  \href{https://huggingface.co/blog/static-embeddings}{Blog post}.
\item
  {[}Cho and Hariharan 2019{]} Jang Hyun Cho and Bharath Hariharan. On
  the Efficacy of Knowledge Distillation. ICCV, 2019.
  \href{https://openaccess.thecvf.com/content_ICCV_2019/html/Cho_On_the_Efficacy_of_Knowledge_Distillation_ICCV_2019_paper.html}{Proceedings}.
\item
  {[}PROD{]} Zhenghao Lin, Yeyun Gong, Xiao Liu, et al.~PROD:
  Progressive Distillation for Dense Retrieval. WWW, 2023.
  \href{https://arxiv.org/abs/2209.13335}{arXiv:2209.13335}.
\item
  {[}Arabzadeh et al.{]} Negar Arabzadeh, Xinyi Yan, and Charles L. A.
  Clarke. Predicting Efficiency/Effectiveness Trade-offs for Dense
  vs.~Sparse Retrieval Strategy Selection. CIKM, 2021.
  \href{https://arxiv.org/abs/2109.10739}{arXiv:2109.10739}.
\item
  {[}QPP{]} Guglielmo Faggioli, Thibault Formal, Stefano Marchesin,
  Stéphane Clinchant, Nicola Ferro, and Benjamin Piwowarski. Query
  Performance Prediction for Neural IR: Are We There Yet? 2023.
  \href{https://arxiv.org/abs/2302.09947}{arXiv:2302.09947}.
\item
  {[}BCT{]} Yantao Shen, Yuanjun Xiong, Wei Xia, and Stefano Soatto.
  Towards Backward-Compatible Representation Learning. CVPR, 2020.
  \href{https://openaccess.thecvf.com/content_CVPR_2020/html/Shen_Towards_Backward-Compatible_Representation_Learning_CVPR_2020_paper.html}{Proceedings}.
\item
  {[}FCT{]} Vivek Ramanujan, Pavan Kumar Anasosalu Vasu, Ali Farhadi,
  Oncel Tuzel, and Hadi Pouransari. Forward Compatible Training for
  Large-Scale Embedding Retrieval Systems. CVPR, 2022.
  \href{https://arxiv.org/abs/2112.02805}{arXiv:2112.02805}.
\item
  {[}OOD-DiskANN{]} Shikhar Jaiswal, Ravishankar Krishnaswamy, Ankit
  Garg, Harsha Vardhan Simhadri, and Sheshansh Agrawal. OOD-DiskANN:
  Efficient and Scalable Graph ANNS for Out-of-Distribution Queries.
  2022. \href{https://arxiv.org/abs/2211.12850}{arXiv:2211.12850}.
\item
  {[}RoarGraph{]} Meng Chen, Kai Zhang, Zhenying He, Yinan Jing, and X.
  Sean Wang. RoarGraph: A Projected Bipartite Graph for Efficient
  Cross-Modal Approximate Nearest Neighbor Search. PVLDB 17(11):
  2735-2749, 2024.
  \href{https://www.vldb.org/pvldb/vol17/p2735-chen.pdf}{Paper}.
\item
  {[}QA-Cos{]} Daehun Nyang. Beyond Hamming: Query-Aware Decoding of
  Binary Cosine Sketches. ICML 2026, PMLR 306: 94126-94143.
  \href{https://proceedings.mlr.press/v306/nyang26a.html}{Proceedings}.
\item
  {[}Qdrant 1.19.1{]} Qdrant contributors. Binary query encoding and
  vector-index search implementation, v1.19.1.
  \href{https://github.com/qdrant/qdrant/tree/v1.19.1}{Tagged source}.
\item
  {[}Stella{]} Stella-en-400M-v5. Model card for
  \texttt{NovaSearch/stella\_en\_400M\_v5}.
  \href{https://huggingface.co/NovaSearch/stella_en_400M_v5}{Model
  page}.
\item
  {[}BEIR{]} BEIR contributors. BEIR retrieval benchmark, datasets, and
  evaluation software.
  \href{https://github.com/beir-cellar/beir}{Repository}.
\end{itemize}

\end{document}
